\subsubsection{Пилот минимизации стоимости при заданной сеточной точности}
\label{sec:costpilot}

Преимущество по ошибке при фиксированных параметрах ещё не определяет,
возникает ли полезная задача выбора метода по времени. В отдельном
ограниченном пилоте~\cite{savedcost60} сравнивались только RK4 и
мультипликативный интеграл (PI) при $N=32,64,128,256$ и допуске
$E_{m,33}\le0.01$. Этот пилот не включал обучение моделей.

\paragraph{Выборка и протокол измерения\\}
Детерминированно отобраны 60 задач исходной обучающей части из 60 разных
семейств и топологий: по три задачи из каждой из 20 ячеек
«сложность~$\times$~режим». Представитель семейства задавался хешем
идентификатора с фиксированным seed, разнообразие коэффициентов
обеспечивалось maximin-отбором по нормированным входным характеристикам.
Ошибки, времена, прогнозы моделей и значения точного решения при выборе
не использовались. Исторические тестовые семейства и топологии были
исключены. Это целевая разнообразная выборка, а не случайная выборка,
по которой можно оценить частоты в генеральной совокупности.

Оба неизменённых решателя получали одни и те же коэффициенты задачи Коши
на $[-1,1]$ и точные начальные условия $f(-1),f'(-1)$.
Для каждого $N$ контрольные узлы $G_{33}$ точно совпадают с узлами
траектории с индексами $kN/32$, $k=0,\ldots,32$; интерполяция не
применялась. Порог включает равенство. Во всех сравнениях используется
именно дискретная ошибка~\eqref{eq:grid33-error}.

На одном компьютере в одном процессе для каждой из 480 пар
«задача~$\times$~конфигурация» выполнены один прогрев и пять измеряемых
повторов: 480 прогревочных и 2400 измеряемых вызовов. Порядок вызовов
детерминированно перемешивался в каждом раунде. Основной показатель
$t_i(m,N)$ --- медиана пяти времён публичного solver-вызова, включающего
вычисление коэффициентов, численные шаги и создание результата.
Загрузка, подготовка точных значений, вычисление ошибки и запись на диск
в этот показатель не включены. Сборщик мусора отключался на измеряемой
фазе, его вызовы выполнялись вне таймера.

Отдельно зарегистрированы импорт (0.208656~с), загрузка и валидация
датасета (2.885793~с), подготовка выбранных входов и точных значений
с вводом-выводом (0.913157~с). Окружение: macOS~26.6.2, arm64,
Python~3.10.11, NumPy~2.2.6, SciPy~1.15.3. Переменные ограничения
внутренних потоков были установлены в~1 до импорта библиотек, однако
\code{threadpoolctl} вернул пустой список. Тепловые и performance-сведения
через \code{pmset} получить не удалось. Поэтому фактическое число
внутренних потоков и отсутствие теплового ограничения частоты независимо
не подтверждены; времена характеризуют эту машину и данный протокол.

\begin{table}[!htbp]
\centering
\caption{Пилот стоимости: точность восьми конфигураций на 60 задачах.
Время --- среднее по задачам от медианы пяти повторов solver;
оно не включает подготовку данных.}
\label{tab:costpilot-configurations}
\small
\begin{tabular}{lrcrc}
\toprule
Метод & $N$ & $E_{m,33}\le0.01$ & Время, мс & $\max E_{m,33}$\\
\midrule
RK4 & 32  & 59/60 & 60.077893  & $1.366209\cdot10^{-2}$\\
RK4 & 64  & 60/60 & 119.962860 & $2.747976\cdot10^{-4}$\\
RK4 & 128 & 60/60 & 239.217672 & $1.699437\cdot10^{-5}$\\
RK4 & 256 & 60/60 & 478.999962 & $1.058317\cdot10^{-6}$\\
PI  & 32  & 0/60  & 15.686458  & 15.916198\\
PI  & 64  & 0/60  & 31.300706  & 8.039300\\
PI  & 128 & 4/60  & 62.703981  & 4.037363\\
PI  & 256 & 7/60  & 125.386115 & 2.022754\\
\bottomrule
\end{tabular}
\end{table}

\paragraph{Oracle и постоянные правила\\}
Oracle для каждой задачи выбирает конфигурацию с минимальной измеренной
медианой среди удовлетворяющих допуску. Лучшая постоянная конфигурация
должна удовлетворять допуску на всех 60 задачах и минимизировать сумму
медиан. Для отделения выбора метода от выбора шага дополнительно
рассмотрен oracle, которому доступны только конфигурации RK4.
Все эти правила определены по уже известным точным ошибкам. Они дают
ретроспективное сравнение и не являются готовыми адаптивными алгоритмами:
стоимость пробных решений и поиска допустимого $N$ в них не учтена.

\begin{table}[!htbp]
\centering
\caption{Стоимость выбора при одинаковом выполнении допуска на 60/60
задачах. Сумма складывается из медиан отдельных solver-вызовов.}
\label{tab:costpilot-selection}
\small
\begin{tabular}{p{0.48\textwidth}rr}
\toprule
Правило & Среднее, мс & Сумма, с\\
\midrule
Лучшая постоянная конфигурация: RK4, $N=64$
 & 119.962860 & 7.197772\\
Oracle по RK4 и PI
 & 62.038960 & 3.722338\\
Oracle только по RK4
 & 62.038960 & 3.722338\\
\bottomrule
\end{tabular}
\end{table}

Во всех 60 случаях oracle выбрал RK4 (100\%); PI не выбран ни разу
(0\%). Для 59 задач достаточно $N=32$, для одной --- $N=64$.
Исключение --- \code{AMV1\_F0186\_S00}: при $N=32$
$E_{\mathrm{RK4},33}=0.0136620881782$, при $N=64$ ---
$0.000210115440205$. Задач, для которых ни одна из восьми конфигураций
не достигла допуска, нет; неизвестных исходов из-за отказов также нет.
PI имеет хотя бы одну допустимую конфигурацию на семи задачах,
но на каждой уступает по медиане подходящей конфигурации RK4.

Среднее относительное сокращение времени oracle относительно постоянного
RK4 при $N=64$ равно $49.1967\%$, а сокращение суммы времён ---
$48.2849\%$. Это разные показатели:
\[
\frac1{60}\sum_{i=1}^{60}
\left(1-\frac{t_i^{\mathrm{oracle}}}{t_i^{\mathrm{const}}}\right),
\qquad
1-\frac{\sum_i t_i^{\mathrm{oracle}}}{\sum_i t_i^{\mathrm{const}}}.
\]
Средняя абсолютная экономия составляет 57.923901~мс на задачу.
Весь выигрыш достигается выбором числа шагов внутри RK4:
добавление PI к RK4-only oracle даёт ровно $0\%$.
Правило минимального допустимого $N$ для RK4 совпало с этим oracle
на всех 60 задачах.

До запуска были зафиксированы критерии GO: не менее 12 из 60 побед
каждого метода и не менее 12 устойчивых побед каждого; не менее 10\%
выигрыша по обоим приведённым показателям относительно лучшей константы
и RK4-only oracle; полное покрытие допуска и не менее 90\% стабильных
профилей. Устойчивая победа требует запаса по медиане не менее 10\%,
неперекрытия интервалов наблюдаемых времён
$\max t_{\rm fast}<\min t_{\rm slow}$ и
$\mathrm{MAD}/\mathrm{median}\le0.20$ у обеих конфигураций.
Это операционные критерии практической полезности, а не статистический
тест значимости. Все 480 профилей удовлетворили ограничению по MAD;
56 методных побед RK4 устойчивы, в четырёх практический запас не
подтверждён. Эти четыре случая не являются победами PI.

\textbf{Результат пилота --- NO-GO для дальнейшего ML выбора между RK4
и PI в данной постановке.} Не получены неоднородность победителей,
устойчивые победы обоих методов и добавочная экономия относительно
выбора $N$ только внутри RK4. Полезность предсказания редкой необходимости
увеличить $N$ отдельно не проверена. Данные пилота не доказывают
обучаемость селектора, переносимость времён, порядок сходимости или
универсальное превосходство RK4. Сохранённая независимая проверка
сверила все ячейки CSV, 2880 ошибок по вложенным узлам и 2400 измерений
времени без новых solver-вызовов~\cite{savedcost60}.

\clearpage
\subsubsection{Граничное продолжение Ремизова на development-выборке}
\label{sec:boundarydev}

Предыдущая диагностика двух экстремальных ответов~\cite{savedboundary}
показала сильное уменьшение ошибки при замене внешнего продолжения
\code{cubic/taper} на \code{cubic/zero}. Для проверки границ этого
наблюдения выполнено одно ограниченное исследование~\cite{savedboundarydev}.
Оно проверяет конкретную реализацию на конечном пространственном окне,
сохраняя математическую формулу метода и исходный производственный
профиль \code{cubic/taper}.

\paragraph{Выборка и контролируемое изменение\\}
Включены все 22 доступные development-задачи с исходной
$E_{\mathrm{R},33}>0.05$ и 30 заранее выбранных контролей с
$E_{\mathrm{R},33}\le0.05$. Проблемная группа охватывает 14 семейств;
контроли взяты из 30 других попарно различных семейств с
детерминированной стратификацией по 20 ячейкам сложности и режима.
Исторические тестовые семейства и топологии исключены, результаты
\code{zero} при отборе не использовались. Три отсутствующих исходных
development-ответа Ремизова не имеют измеренной $E_{\mathrm{R},33}$,
поэтому не входят в эти 52 задачи и не заменяются условной ошибкой.
Прежняя задача \code{F0180/S04} принадлежит Test и в новую оценку не
включалась; в исторической диагностике её ошибка уменьшилась с
19134.405 до 0.0704519 и осталась выше допуска 0.05~\cite{savedboundary}.

Для всех 52 задач повторно использованы исходные ответы \code{taper};
два полностью совместимых ответа \code{zero} также взяты из сохранённых
результатов. Выполнено ровно 50 новых вызовов \code{zero} и ни одного
нового вызова \code{taper}. В сравнении изменён только параметр внешнего
продолжения. Сохранены $n=128$, 513 временных узлов с квадратичным
размещением, $\Delta x=0.0125$, дополнение окна на 12 с каждой стороны
интервала $[-1,1]$, то есть окно $[-13,13]$ из 2081 узла,
и $\code{tail\_eps}=10^{-8}$. Эффективное $T$ совпадает для каждого
идентификатора. Совпадение коэффициентов, точного решения,
пространственных и временных сеток, квадратурных весов и $G_{33}$
проверено по исходным значениям и контрольным суммам. Другие продолжения,
MC и ML в этой серии не исследовались.

Порог точности равен $0.05$. Катастрофическим выбросом в протоколе
условно названо $E_{\mathrm{R},33}>1$; это определение исхода,
а не утверждение о математической неприменимости метода.
Для разности $\Delta E=E^{\mathrm{zero}}-E^{\mathrm{taper}}$
заранее задан допуск численного сравнения
\[
\delta_i=10^{-10}\max\{1,E_i^{\mathrm{taper}}\}.
\]
Уменьшение означает $\Delta E<-\delta_i$, совпадение ---
$|\Delta E|\le\delta_i$, регрессия по этому критерию ---
$\Delta E>\delta_i$. Этот допуск не является ни границей полной ошибки,
ни порогом практической существенности.

\begin{table}[!htbp]
\centering
\caption{Результат замены \code{taper} на \code{zero}.
Уменьшения, совпадения и регрессии определены с допуском $\delta_i$.
Столбцы с парой значений относятся к исходному и новому продолжению.}
\label{tab:boundarydev-groups}
\resizebox{\textwidth}{!}{%
\begin{tabular}{lrrrrrr}
\toprule
Группа & $n$ & \parbox{23mm}{\centering $E_{33}\le0.05$, до/после} & Снижений & Равенств
 & Регрессий & \parbox{20mm}{\centering $E_{33}>1$, до/после}\\
\midrule
Проблемные & 22 & 0/1 & 9 & 10 & 3 & 3/0\\
Контроли & 30 & 30/30 & 2 & 28 & 0 & 0/0\\
Все & 52 & 30/31 & 11 & 38 & 3 & 3/0\\
\bottomrule
\end{tabular}}
\end{table}

\paragraph{Устранённые выбросы и оставшиеся ошибки\\}
Все три исходных выброса выше единицы опустились ниже единицы
(таблица~\ref{tab:boundarydev-extremes}). Это подтверждает влияние
граничного продолжения на данные выбросы при фиксированных остальных
параметрах. Однако только одна из 22 проблемных задач перешла ниже 0.05;
21 задача осталась вне допуска. Максимум по 52 задачам уменьшился
с 19704.5019109 до 0.596851548154. У всех 30 контролей допуск сохранён;
новых выбросов выше единицы и неуспешных ответов \code{zero} нет.

\begin{table}[!htbp]
\centering
\caption{Три исходных катастрофических выброса development-выборки.
Приведены итоговые ошибки на $G_{33}$, а не промежуточные состояния.}
\label{tab:boundarydev-extremes}
\small
\begin{tabular}{lrrc}
\toprule
Задача & $E_{33}$, \code{taper} & $E_{33}$, \code{zero}
 & Допуск после\\
\midrule
\code{F0141/S00} & 19704.5019109 & 0.0187192785409 & да\\
\code{F0184/S01} & 444.589254639 & 0.0535680732159 & нет\\
\code{F0186/S02} & 2.24644415774 & 0.254377736872 & нет\\
\bottomrule
\end{tabular}
\end{table}

Положительная неокруглённая разность наблюдается у 19 задач, включая
увеличения меньше $\delta_i$. Только три разности превышают заранее
заданный допуск сравнения (таблица~\ref{tab:boundarydev-regressions}).
Все три задачи уже находились выше 0.05; новых потерь допуска нет.
Приращения порядка $10^{-9}$ не свидетельствуют об установленном
практически значимом ухудшении.

\begin{table}[!htbp]
\centering
\caption{Все три регрессии, превышающие $\delta_i$.
Относительное увеличение приведено как доля, а не процент.}
\label{tab:boundarydev-regressions}
\small
\begin{tabular}{lcc}
\toprule
Задача & $\Delta E_{33}$ & $\Delta E_{33}/E_{33}^{\mathrm{taper}}$\\
\midrule
\code{F0186/S01} & $6.098373472873675\cdot10^{-10}$ & $1.9713427\cdot10^{-9}$\\
\code{F0189/S02} & $1.312633113315087\cdot10^{-9}$ & $5.4572689\cdot10^{-9}$\\
\code{F0189/S03} & $6.146191888767305\cdot10^{-10}$ & $7.2482444\cdot10^{-9}$\\
\bottomrule
\end{tabular}
\end{table}

\paragraph{Пространственная аппроксимация и теоретические ограничения\\}
В \code{zero} неизменённый внутренний кубический сплайн используется
внутри окна, а запросам строго за окном присваивается ноль.
Граничный узел при этом сохраняет своё вычисленное значение. Если оно
ненулевое, продолжение разрывно и не принадлежит пространствам
непрерывных ограниченных функций $C_b$ или $UC_b$, когда они требуются
в теоретическом обосновании. Это не отдельный решатель корректно
поставленной краевой задачи с условием Дирихле: граничные узловые
значения принудительно не обнуляются. Даже при нулевых концах гладкость
продолжения не гарантирована. Например, сплайн через узлы
$x=(0,1,2,3)$ со значениями $(0,1,1,0)$ равен
$p(x)=3x/2-x^2/2$; его производные на концах равны $3/2$ и $-3/2$,
а производная внешнего нуля равна нулю.

Независимые сохранённые проверки~\cite{savedboundarydev} показывают
также отсутствие гарантии положительности и контрактности
\emph{дискретной кубической реконструкции в узловой максимальной норме}.
При $x=(0,1,2,3)$, $a=1/4$, $b=4$, $c=0$, $\tau=1/16$
и \code{boundary=zero} неизменённая реализация одного применения
$\Sop(\tau)$ даёт в узле $x=1$ значение $67/64>1$ для данных
$u=(0,1,1,0)$ с узловой нормой~1. Для неотрицательных узловых данных
$u=(1,0,0,0)$ результат в том же узле равен $-1/64$.
Это контрпримеры для операции «кубическая реконструкция плюс сдвиги»
на векторе узловых значений, а не для точного оператора $\Sop$ на
непрерывных функциях: сама реконструкция может превышать диапазон
узловых данных. Формальные веса $1/4,1/4,1/2$ не дают автоматически
оценку нормы полученного сеточного оператора.

Исходная правая часть имеет носитель внутри окна, но повторные сдвиги
распространяют промежуточные состояния. Поэтому начальная компактность
носителя не исключает последующего влияния границ. Сохранённая величина
\code{boundary\_state\_max\_abs} характеризует эти промежуточные
состояния и не равна итоговой $E_{\mathrm{R},33}$.
Формула $\Sop$, коэффициенты и их глобальные сертификаты сохранены;
пространственная устойчивость, погрешность усечения окна и выполнение
всех предпосылок глобальной теоремы для этой конечномерной реализации
не сертифицированы. Уточнения $\Delta x$, окна, $T$, квадратуры и $n$
в данном сравнении не выполнялись, поэтому причинные доли оставшихся
источников полной ошибки не установлены.

До получения новых ответов GO требовал пригодности всех результатов,
отсутствия регрессий по $\delta_i$, потерь допуска и новых выбросов выше
единицы, устранения всех прежних выбросов выше единицы, увеличения числа
точных ответов и хотя бы одного уменьшения сверх $\delta_i$.
Не выполнено только условие отсутствия регрессий.
\textbf{Формальный итог --- NO-GO для принятия \code{cubic/zero}
как нового экспериментального профиля без регрессий по этому протоколу.}
Критерий не ослаблялся после наблюдений. Этот итог чувствителен к
строгости численного порога и не означает доказанной практической
вредности \code{zero}; сильное уменьшение трёх выбросов остаётся
подтверждённым результатом. Вместе с тем \code{zero} не устраняет
все крупные ошибки и не обеспечивает автоматически теоретическую
корректность пространственной аппроксимации. Исходный производственный
профиль и основные массовые таблицы не заменялись.
Независимая проверка сверила 1664 ячейки таблицы, все 52 ошибки каждого
профиля, два повторно использованных ответа и журнал
50 новых вызовов; проверка не решала ОДУ повторно~\cite{savedboundarydev}.
